\documentclass[11pt,a4paper]{article}
\usepackage[T1]{fontenc}
\usepackage[utf8]{inputenc}
\usepackage{lmodern}
\usepackage[czech]{babel}
\usepackage[margin=2cm]{geometry}
\usepackage{amsmath,amssymb,mathtools}
\usepackage{xcolor,cancel,enumitem}
\usepackage{tikz,pgfplots}
\pgfplotsset{compat=1.18}
\usetikzlibrary{arrows.meta,calc,patterns}
\pagestyle{empty}
\setlength{\parindent}{0pt}
\setlength{\parskip}{4pt}
\newcommand{\vysl}[1]{\boxed{\color{red!75!black}#1}}
\newcommand{\pozn}[1]{{\color{red!75!black}#1}}

% osa s bodem a šipkou ukazující, odkud se blížíme
\newcommand{\priblizeni}[2]{%
\begin{tikzpicture}[baseline=-3pt]
\draw (-2,0) -- (2,0);
\draw (0,0.1) -- (0,-0.1) node[below] {$#1$};
\draw[red!75!black, -{Stealth}] (1.6,-0.35) -- (0.3,-0.35);
\node[red!75!black, below] at (1.2,-0.4) {\small $#2$};
\end{tikzpicture}}

\begin{document}

\textbf{c)}
\[
\lim_{x\to-3^+} \frac{2-x}{2x^2+5x-3}
= \left\|\frac{5}{0}\right\|
= \lim_{x\to-3^+} \frac{2-x}{2\left(x-\frac12\right)(x+3)} =
\]

Rozklad jmenovatele $2x^2+5x-3$:
\begin{align*}
D &= b^2-4ac\\
D &= 25-4\cdot 2\cdot(-3)\\
D &= 25+24 = 49\\
x_{1,2} &= \frac{-b\pm\sqrt{D}}{2a} = \frac{-5\pm 7}{4} =
\begin{cases} \frac{2}{4} = \frac12\\[3pt] -\frac{12}{4} = -3 \end{cases}
\end{align*}

\[
= \frac{5}{2\cdot(-3{,}5)\cdot(0{,}1)} = \frac{5}{0^-} = \vysl{-\infty}
\qquad
\priblizeni{-3}{-2{,}9}
\]

\pozn{Kde mi nevychází nula, tam dosadím rovnou $x=-3$,
kde vychází nula, tam dám hodnotu $-2{,}9$!
Abych věděla, zda je to číslo blízko nuly KLADNÉ či ZÁPORNÉ.}

\bigskip
\textbf{d)}
\[
\lim_{x\to4^+} \ln\frac{3x+1}{x^2-4x}
= \ln\left(\lim_{x\to4^+}\frac{3x+1}{x^2-4x}\right)
= \ln(+\infty) = +\infty
\]
(funkce $\ln x$ roste pro $x\to+\infty$ nade všechny meze)

Vnitřní limita:
\[
\lim_{x\to4^+} \frac{3x+1}{x^2-4x}
= \left\|\frac{13}{0}\right\|
= \lim_{x\to4^+} \frac{3x+1}{x(x-4)}
= \frac{13}{4\cdot(0{,}1)}
= \frac{13}{4\cdot0^+}
= \frac{13}{0^+}
= \pozn{+\infty}
\qquad
\priblizeni{4}{4{,}1}
\]

\bigskip
\textbf{e)}
\[
\lim_{x\to-\infty} \left(x^3-x+2\right)
= \lim_{x\to-\infty} x^3
= (-\infty)^3 = \vysl{-\infty}
\]

\end{document}
